\subsection{Parte intera, densità e intervalli}
\label{sec:esercizi-fondamenti-parte-intera}

\begin{exercise}[Parte intera inferiore e superiore di numeri negativi]
\label{ex:fondamenti-floor-ceiling-negativi}
\difficoltaesercizio{0}
Calcolare $\lfloor-7/3\rfloor$, $\lceil-7/3\rceil$,
$\lfloor5/2\rfloor$ e $\lceil5/2\rceil$, motivando ogni
risultato mediante le rispettive definizioni.
\rimandoteoria{def:floor,def:ceiling}
\end{exercise}
\begin{proof}[Soluzione dell'esercizio]
Poiché
\[
  -3<-\frac73<-2,
  \qquad 2<\frac52<3,
\]
per le definizioni di floor e ceiling si ottiene
\[
  \lfloor-7/3\rfloor=-3,\quad
  \lceil-7/3\rceil=-2,\quad
  \lfloor5/2\rfloor=2,\quad
  \lceil5/2\rceil=3.
\]
\end{proof}

\begin{exercise}[Equazione con la parte intera inferiore]
\label{ex:fondamenti-equazione-parte-intera}
\difficoltaesercizio{0}
Determinare tutti i numeri reali $x$ tali che
\[
  \lfloor2x-1\rfloor=3.
\]
Scrivere l'insieme soluzione come intervallo.
\rimandoteoria{def:floor}
\end{exercise}
\begin{proof}[Soluzione dell'esercizio]
Per la definizione di floor,
\[
  \lfloor2x-1\rfloor=3
  \quad\Longleftrightarrow\quad
  3\le2x-1<4
  \quad\Longleftrightarrow\quad
  2\le x<\frac52.
\]
L'insieme soluzione è quindi $[2,5/2)$.
\end{proof}

\begin{exercise}[Approssimazione razionale di un irrazionale]
\label{ex:fondamenti-approssimazione-razionale-irrazionale}
\difficoltaesercizio{1}
Trovare esplicitamente un numero $q\in\Q$ tale che
\[
  0<\sqrt2-q<\frac1{20}.
\]
Verificare entrambe le disuguaglianze senza ricorrere
ad approssimazioni decimali non giustificate.
\rimandoteoria{thm:densita-razionali}
\end{exercise}
\begin{proof}[Soluzione dell'esercizio]
Si sceglie $q=7/5\in\Q$. Poiché
\[
  \left(\frac75\right)^2=\frac{49}{25}<2
  <\frac{841}{400}=\left(\frac{29}{20}\right)^2,
\]
e i numeri considerati sono positivi, segue
\[
  \frac75<\sqrt2<\frac{29}{20}
  =\frac75+\frac1{20}.
\]
Pertanto $0<\sqrt2-q<1/20$.
\end{proof}

\begin{exercise}[Parte intera di numeri opposti]
\label{ex:fondamenti-parte-intera-numeri-opposti}
\difficoltaesercizio{2}
Dimostrare che, per ogni $x\in\R$,
\[
  \lfloor x\rfloor+\lfloor-x\rfloor\in\{-1,0\}.
\]
Determinare esattamente per quali $x$ si verifica
ciascuno dei due valori.
\rimandoteoria{def:floor}
\end{exercise}
\begin{proof}[Soluzione dell'esercizio]
Si ponga $n:=\lfloor x\rfloor$. Per definizione,
\[
  n\le x<n+1,
  \qquad -n-1<-x\le-n.
\]
Se $x\in\Z$, si ha $x=n$ e $\lfloor-x\rfloor=-n$;
pertanto la somma è nulla.

Se $x\notin\Z$, risulta $n<x<n+1$ e quindi
$-n-1<-x<-n$. Ne segue $\lfloor-x\rfloor=-n-1$,
per cui la somma vale $-1$. In conclusione,
\[
  \lfloor x\rfloor+\lfloor-x\rfloor=
  \begin{cases}
    0,&x\in\Z,\\
    -1,&x\in\R\setminus\Z.
  \end{cases}
\]
\end{proof}
